% !TEX program = pdflatex
\documentclass[11pt]{article}

% ── Packages ──────────────────────────────────────────────────────────
\usepackage[margin=1in]{geometry}
\usepackage{amsmath,amssymb,amsthm}
\usepackage{booktabs}
\usepackage{graphicx}
\usepackage{hyperref}
\usepackage{cleveref}
\usepackage{xcolor}
\usepackage{caption}
\usepackage{subcaption}
\usepackage{multirow}
\usepackage[numbers]{natbib}
\bibliographystyle{plainnat}

\hypersetup{colorlinks=true,linkcolor=blue!60!black,citecolor=green!50!black,urlcolor=blue!70!black}

% ── Macros ────────────────────────────────────────────────────────────
\newcommand{\Einf}{E_\infty}
\newcommand{\Reval}{R_{\mathrm{eval}}}
\newcommand{\Rtrain}{R_{\mathrm{train}}}
\newcommand{\relLtwo}{\varepsilon_{\mathrm{rel}}}
\DeclareMathOperator*{\argmin}{arg\,min}

% ══════════════════════════════════════════════════════════════════════
\title{%
  Scaling Laws Reveal Resolution as a First-Class Variable \\
  in Operator Learning%
}
\author{Lucas Perrier}
\date{\today \quad (Working Draft)}

\begin{document}
\maketitle

% ══════════════════════════════════════════════════════════════════════
\begin{abstract}
Neural operators promise mesh-independent learning of PDE solution maps, yet
practitioners lack quantitative guidance on how test error scales with dataset
size, model capacity, and discretization resolution. We conduct a
comprehensive scaling study across three PDE benchmarks (Burgers, Darcy,
diffusion), three architectures (MLP baseline, DeepONet, FNO), and 5{,}040
training runs spanning two orders of magnitude in each scaling axis. We fit
the \emph{three-axis scaling law}
$E(N,D,R) = \Einf + a\,N^{-\alpha} + b\,D^{-\beta} + c\,R^{-\gamma}$
and find that (i)~data-scaling exponents $\alpha$ are strongly
task-dependent, ranging from $\approx 0.4$ (Burgers, Darcy) to
$\approx 1.1$ (diffusion); (ii)~resolution acts as a meaningful
scaling axis primarily for advection-dominated problems (Burgers $\gamma
\approx 0.5$--$1.3$), while diffusion-dominated problems saturate
resolution benefits rapidly ($\gamma \to 5$, the fitting bound);
(iii)~FNO achieves the lowest absolute errors across all tasks but
derives the least relative benefit from additional resolution on
Burgers ($\gamma_{\mathrm{FNO}} \approx 0.77$) compared with DeepONet
($\gamma_{\mathrm{DeepONet}} \approx 1.28$); and (iv)~cross-resolution
transfer matrices reveal that FNO's spectral inductive bias enables
strong generalization from high-$\Rtrain$ to unseen resolutions,
while DeepONet transfers more uniformly but with higher absolute error.
These findings establish resolution as a first-class variable for
operator-learning scaling laws and provide actionable guidance for
practitioners budgeting compute across data collection, model sizing,
and grid refinement.
\end{abstract}

% ══════════════════════════════════════════════════════════════════════
\section{Introduction}
\label{sec:intro}

Empirical scaling laws have become a central tool for understanding and
predicting the performance of large-scale machine learning
systems~\citep{kaplan2020scaling,hoffmann2022training}. In natural
language processing and computer vision, power-law relationships between
test loss and dataset size, parameter count, and compute budget guide
resource allocation decisions worth millions of dollars.

Scientific machine learning (SciML) faces an analogous resource
allocation problem with a twist: the practitioner must additionally
choose the \emph{discretization resolution} at which training data are
generated. Solving PDEs at higher resolution is expensive---often the
dominant cost---yet the relationship between resolution and downstream
operator-learning error is poorly characterized.

Neural operators~\citep{lu2021learning,li2021fourier} learn mappings
between infinite-dimensional function spaces and are often marketed as
``resolution-invariant.'' Two prominent architectures---DeepONet
\citep{lu2021learning} and the Fourier Neural Operator (FNO)
\citep{li2021fourier}---use fundamentally different strategies to
handle varying discretizations. DeepONet factorizes the operator into
branch (input encoding) and trunk (output basis) networks, making it
architecturally independent of the input grid. FNO parameterizes
integral kernels in Fourier space, enabling it to operate at any
resolution that supports the retained frequency modes.

Despite these architectural innovations, it remains unclear \emph{when}
resolution invariance translates into practical benefits. Does doubling
the training resolution halve the error? Does it matter more than
doubling the dataset? How does the answer depend on the PDE?

\paragraph{Contributions.}
\begin{enumerate}
  \item We propose the \textbf{three-axis scaling law}
    $E(N,D,R) = \Einf + a\,N^{-\alpha} + b\,D^{-\beta} + c\,R^{-\gamma}$,
    which adds discretization resolution $R$ as a first-class scaling
    axis alongside dataset size $N$ and model capacity $D$.
  \item We report \textbf{5{,}040 training runs} and
    \textbf{13{,}440 cross-resolution evaluations} across three PDE
    families, three architectures, and seven capacity levels, providing
    the most comprehensive operator-learning scaling study to date.
  \item We find that the resolution exponent $\gamma$ is
    \textbf{strongly task-dependent}: large for advection-dominated PDEs
    (Burgers, $\gamma \approx 0.5$--$1.3$) and negligible for
    diffusion-dominated PDEs ($\gamma \to 5$, i.e., rapid saturation).
  \item We introduce \textbf{cross-resolution transfer matrices} that
    reveal how operator-learning models generalize across resolutions,
    showing that FNO's spectral bias enables strong high-to-low
    resolution transfer on Burgers while DeepONet transfers more
    uniformly.
\end{enumerate}


% ══════════════════════════════════════════════════════════════════════
\section{Related Work}
\label{sec:related}

\paragraph{Neural scaling laws.}
\citet{kaplan2020scaling} and \citet{hoffmann2022training} established
power-law scaling of cross-entropy loss with dataset tokens, parameter
count, and compute for autoregressive language models.
\citet{hestness2017deep} showed similar trends across vision and speech.
In scientific ML, \citet{perrier2025scaling} introduced scaling-law
diagnostics for physics-informed approaches (PINNs, HNNs) in
low-dimensional settings, finding that physics priors shift the error
floor rather than the scaling exponent. Our work extends this program to
infinite-dimensional operator learning and elevates resolution to a
central variable.

\paragraph{Operator learning.}
DeepONet~\citep{lu2021learning} builds on the universal approximation
theorem for operators~\citep{chen1995universal}, factorizing the learned
operator into branch and trunk components. The Fourier Neural Operator
(FNO)~\citep{li2021fourier} parameterizes integral kernels via
truncated Fourier series, achieving state-of-the-art accuracy on
several PDE benchmarks. Recent extensions include U-shaped
architectures~\citep{wen2022ufno}, attention
mechanisms~\citep{hao2023gnot}, and physics-informed
variants~\citep{wang2021learning}.

\paragraph{Resolution and discretization in operator learning.}
A key selling point of neural operators is their ability to generalize
across resolutions. \citet{li2021fourier} demonstrated FNO's
zero-shot super-resolution capability. However, systematic studies of
how \emph{training} resolution affects accuracy---and its interaction
with dataset size and model capacity---are lacking. Our work fills this
gap with controlled experiments isolating the resolution axis.


% ══════════════════════════════════════════════════════════════════════
\section{Scaling Law Framework}
\label{sec:framework}

\subsection{Problem Setting}
We consider the supervised learning of solution operators
$\mathcal{G}^\dagger : \mathcal{U} \to \mathcal{V}$, where
$\mathcal{U}$ and $\mathcal{V}$ are Banach spaces of functions on
some domain $\Omega \subset \mathbb{R}^d$. Given $N$ input--output
pairs $\{(u_i, v_i)\}_{i=1}^N$ with
$v_i = \mathcal{G}^\dagger(u_i)$, a model $\mathcal{G}_\theta$ with
$D = |\theta|$ parameters is trained on data discretized at spatial
resolution $R$ (grid points per spatial dimension).

\subsection{Three-Axis Scaling Law}
We posit that the expected test error follows:
\begin{equation}
  \label{eq:3d-law}
  E(N, D, R) \;=\; \Einf \;+\;
    a\,N^{-\alpha} \;+\; b\,D^{-\beta} \;+\; c\,R^{-\gamma},
\end{equation}
where $\Einf$ is the irreducible error floor (approximation error from
architecture + task complexity), and $\alpha$, $\beta$, $\gamma > 0$
are the data, capacity, and resolution scaling exponents. The additive
decomposition assumes approximate independence of the three error
sources in the power-law regime.

\paragraph{Interpretation of exponents.}
\begin{itemize}
  \item $\alpha$ governs the \emph{statistical efficiency}: how
    quickly additional training data reduces error. Large $\alpha$
    indicates the model learns efficiently from data.
  \item $\beta$ governs the \emph{capacity efficiency}: how
    well the architecture converts additional parameters into accuracy.
    Small $\beta$ suggests the model is already well-matched to the
    task complexity.
  \item $\gamma$ governs the \emph{resolution efficiency}: the
    marginal value of finer discretization. Large $\gamma$ indicates
    rapid saturation (resolution is not a bottleneck); small $\gamma$
    indicates that higher resolution provides lasting benefits.
\end{itemize}

\subsection{Fitting Procedure}
For marginal scaling laws (e.g., $E$ vs.\ $N$ at fixed $D, R$), we fit
$E(X) = \Einf + a\,X^{-\alpha}$ via nonlinear least squares
(Levenberg--Marquardt) with bounds $\Einf \geq 0$,
$a > 0$, $0.01 \leq \alpha \leq 5$. For the full 3D law
\eqref{eq:3d-law}, we fit all seven parameters jointly. Confidence
intervals are obtained via 200 bootstrap resamples over the grouped
data points.


% ══════════════════════════════════════════════════════════════════════
\section{Experimental Setup}
\label{sec:experiments}

\subsection{PDE Benchmarks}
We study three PDE families that span a range of physical behaviors:

\begin{table}[h]
\centering
\caption{PDE benchmark tasks. All tasks map an input function to an
output function on the same domain, discretized at resolution $R$.}
\label{tab:tasks}
\begin{tabular}{@{}llll@{}}
\toprule
\textbf{Task} & \textbf{PDE} & \textbf{Operator} & \textbf{Character} \\
\midrule
Burgers & $u_t + u u_x = \nu u_{xx}$ &
  $u_0(x) \mapsto u(x, T)$ & Advection-dominated \\
Darcy & $-\nabla \cdot (a \nabla u) = f$ &
  $a(x) \mapsto u(x)$ & Elliptic, multiscale \\
Diffusion & $u_t = \nu u_{xx}$ &
  $u_0(x) \mapsto u(x, T)$ & Diffusion-dominated \\
\bottomrule
\end{tabular}
\end{table}

Burgers uses a spectral solver with integrating-factor RK4 and
$\frac{2}{3}$-rule dealiasing. Darcy uses a spectral Galerkin method.
Diffusion uses an exact spectral solution. All solvers are
resolution-parameterized, generating data at arbitrary $R$.

\subsection{Architectures}

\begin{table}[h]
\centering
\caption{Model architectures. The capacity grid spans 7 levels from
\textit{tiny} ($\sim$5K params) to \textit{xlarge} ($\sim$500K params).
MLP parameter counts scale with input resolution $R$; DeepONet and
FNO counts are resolution-independent.}
\label{tab:models}
\begin{tabular}{@{}llp{6cm}@{}}
\toprule
\textbf{Model} & \textbf{Type} & \textbf{Resolution handling} \\
\midrule
MLP & Discretization-tied &
  Flattened input of dimension $R$; parameter count $\propto R$ \\
DeepONet & Mesh-free &
  Branch encodes input at training $R$; trunk evaluates at arbitrary
  query points \\
FNO & Spectral &
  Global convolutions in Fourier space; truncation to $k_{\max}$ modes
  enables resolution transfer \\
\bottomrule
\end{tabular}
\end{table}

\subsection{Scaling Grid}
We sweep over five axes:
\begin{itemize}
  \item \textbf{Dataset size} $N \in \{50, 100, 500, 1000, 5000\}$
  \item \textbf{Resolution} $R \in \{32, 64, 128, 256\}$
  \item \textbf{Model capacity}: 7 levels (\textit{tiny} to
    \textit{xlarge})
  \item \textbf{Data seeds}: $\{11, 22\}$
  \item \textbf{Training seeds}: $\{101, 202\}$
\end{itemize}
This yields $3 \times 7 \times 5 \times 4 \times 2 \times 2 = 1{,}680$
runs per task and $5{,}040$ total training runs across three tasks. Each
model is trained for up to 5{,}000 epochs with early stopping (patience
200) using Adam ($\text{lr} = 10^{-3}$, weight decay $10^{-5}$).

\subsection{Cross-Resolution Evaluation}
After training, each model is evaluated at all resolutions
$\Reval \in \{32, 64, 128, 256, 512\}$ except its training resolution,
yielding 13{,}440 additional evaluation metrics. For DeepONet, branch
inputs at non-training resolutions are linearly interpolated to
$\Rtrain$ before forward pass. FNO natively handles arbitrary input
resolutions via Fourier truncation/zero-padding.


% ══════════════════════════════════════════════════════════════════════
\section{Results}
\label{sec:results}

All 5{,}040 training runs and 13{,}440 cross-resolution evaluations
completed with zero divergence across all model--task combinations.
We report relative $L^2$ error
$\relLtwo = \|v_{\mathrm{pred}} - v_{\mathrm{true}}\|_2 /
             \|v_{\mathrm{true}}\|_2$
as the primary metric.

% ── 5.1: Data scaling ────────────────────────────────────────────────
\subsection{Data Scaling}
\label{sec:data-scaling}

\Cref{fig:data-scaling} shows $\relLtwo$ vs.\ $N$ at medium capacity
and $R = 128$ for each task. The data-scaling exponent $\alpha$ varies
strongly across tasks:

\begin{table}[h]
\centering
\caption{Data-scaling exponents $\alpha$ from the full 3D fits
$E(N,D,R)$, with 95\% bootstrap confidence intervals.}
\label{tab:alpha}
\begin{tabular}{@{}lccc@{}}
\toprule
& \textbf{MLP} & \textbf{DeepONet} & \textbf{FNO} \\
\midrule
Burgers & $0.375\;[0.30, 0.46]$ & $0.632\;[0.44, 0.82]$
        & $0.771\;[0.65, 0.90]$ \\
Darcy   & $0.418\;[0.37, 0.48]$ & $0.459\;[0.34, 0.58]$
        & $0.426\;[0.33, 0.51]$ \\
Diffusion & $1.063\;[0.95, 1.22]$ & $1.273\;[1.00, 1.58]$
          & $1.092\;[0.92, 1.25]$ \\
\bottomrule
\end{tabular}
\end{table}

\paragraph{Key finding: task dependence dominates model dependence.}
The data-scaling exponent is primarily determined by the PDE, not the
architecture. Diffusion ($\alpha \approx 1.1$) shows steep data
scaling---models learn this smooth operator efficiently---while Burgers
($\alpha \approx 0.4$--$0.8$) and Darcy ($\alpha \approx 0.4$--$0.5$)
are harder. For Burgers, FNO benefits most from additional data
($\alpha = 0.77$), likely because its spectral parameterization aligns
with the Fourier-based solver. For diffusion, the MLP achieves an 8.8$\times$
improvement from $N = 50$ to $N = 5{,}000$, the largest data-scaling
gain of any model--task pair.

\begin{figure}[t]
\centering
\includegraphics[width=\textwidth]{../figures/fig_data_scaling_per_task.png}
\caption{Data scaling: relative $L^2$ error vs.\ $N$ at medium
capacity and $R = 128$. Dashed lines are power-law fits
$\Einf + a\,N^{-\alpha}$. Diffusion shows the steepest decay;
Darcy shows the flattest.}
\label{fig:data-scaling}
\end{figure}

% ── 5.2: Resolution scaling ─────────────────────────────────────────
\subsection{Resolution Scaling}
\label{sec:resolution-scaling}

\Cref{tab:gamma} reports the resolution exponent $\gamma$ from the 3D fits.

\begin{table}[h]
\centering
\caption{Resolution-scaling exponents $\gamma$ from the full 3D fits,
with 95\% bootstrap CIs. Values capped at the fitting bound of 5.0
indicate negligible resolution dependence.}
\label{tab:gamma}
\begin{tabular}{@{}lccc@{}}
\toprule
& \textbf{MLP} & \textbf{DeepONet} & \textbf{FNO} \\
\midrule
Burgers & $0.522\;[0.38, 0.73]$ & $1.280\;[1.06, 1.50]$
        & $0.768\;[0.62, 0.95]$ \\
Darcy   & $2.890\;[0.24, 5.00]$ & $3.423\;[0.26, 5.00]$
        & $2.589\;[0.61, 5.00]$ \\
Diffusion & $4.860\;[4.70, 5.00]$ & $5.000\;[5.00, 5.00]$
          & $4.632\;[0.05, 5.00]$ \\
\bottomrule
\end{tabular}
\end{table}

\paragraph{Key finding: resolution matters most for advection.}
For Burgers, which features steep gradients and shock-like structures,
all models benefit from higher resolution. DeepONet shows the largest
$\gamma$ ($1.28$), suggesting its branch--trunk factorization is
especially sensitive to input discretization quality. FNO benefits
moderately ($\gamma = 0.77$), consistent with its built-in spectral
smoothing. At medium capacity and $N = 1{,}000$, increasing $R$ from 32
to 256 yields a 1.7$\times$ error reduction for FNO and 1.6$\times$ for
DeepONet.

For diffusion, all $\gamma$ values hit or approach the fitting bound of
5.0, indicating that even $R = 32$ captures the smooth solution
adequately. The MLP---despite being resolution-tied---shows similarly
negligible resolution sensitivity, confirming that the smoothness of the
operator, not the architecture, determines resolution benefit.

Darcy falls between: wide CIs reflect weak but nonzero resolution
dependence, consistent with the elliptic PDE's moderate spatial
complexity.

\begin{figure}[t]
\centering
\includegraphics[width=\textwidth]{../figures/fig_resolution_scaling_per_task.png}
\caption{Resolution scaling: $\relLtwo$ vs.\ $R$ at medium capacity
and $N = 1{,}000$. Burgers shows clear power-law decay; diffusion is
essentially flat.}
\label{fig:resolution-scaling}
\end{figure}

% ── 5.3: Capacity scaling ───────────────────────────────────────────
\subsection{Capacity Scaling}
\label{sec:capacity-scaling}

\begin{table}[h]
\centering
\caption{Capacity-scaling exponents $\beta$ from the full 3D fits.
Small $\beta$ indicates the model is already well-sized for the
task.}
\label{tab:beta}
\begin{tabular}{@{}lccc@{}}
\toprule
& \textbf{MLP} & \textbf{DeepONet} & \textbf{FNO} \\
\midrule
Burgers & $0.130\;[0.07, 0.27]$ & $0.054\;[0.04, 0.07]$
        & $0.572\;[0.35, 3.91]$ \\
Darcy   & $0.438\;[0.15, 4.05]$ & $0.052\;[0.04, 5.00]$
        & $3.752\;[0.49, 4.06]$ \\
Diffusion & $1.256\;[0.65, 2.50]$ & $0.115\;[0.10, 0.13]$
          & $4.303\;[1.70, 4.32]$ \\
\bottomrule
\end{tabular}
\end{table}

\paragraph{Key finding: DeepONet shows consistently low capacity scaling.}
DeepONet's $\beta$ is small across all tasks ($0.05$--$0.12$),
indicating that even the smallest capacity levels perform comparably to
the largest. This suggests that the branch--trunk architecture is
already a good structural match for these operators, and additional
parameters do not help substantially. FNO, by contrast, shows
high $\beta$ for Darcy and diffusion, indicating it can effectively
leverage additional modes and width.

\begin{figure}[t]
\centering
\includegraphics[width=\textwidth]{../figures/fig_capacity_scaling_per_task.png}
\caption{Capacity scaling: $\relLtwo$ vs.\ parameter count $D$ at
$N = 1{,}000$ and $R = 128$.}
\label{fig:capacity-scaling}
\end{figure}

% ── 5.4: Full 3D fits ───────────────────────────────────────────────
\subsection{Full Three-Dimensional Fits}
\label{sec:3d-fits}

\Cref{tab:3d-r2} summarizes the goodness-of-fit ($R^2$) for the joint
law \eqref{eq:3d-law} across all model--task combinations.

\begin{table}[h]
\centering
\caption{$R^2$ values for the full 3D scaling law
$E(N,D,R) = \Einf + a N^{-\alpha} + b D^{-\beta} + c R^{-\gamma}$.
All fits use 140 grouped data points per model per task.}
\label{tab:3d-r2}
\begin{tabular}{@{}lccc@{}}
\toprule
& \textbf{MLP} & \textbf{DeepONet} & \textbf{FNO} \\
\midrule
Burgers   & 0.935 & 0.915 & 0.871 \\
Darcy     & 0.886 & 0.788 & 0.880 \\
Diffusion & 0.951 & 0.842 & 0.893 \\
\bottomrule
\end{tabular}
\end{table}

The additive three-axis law provides an excellent fit ($R^2 > 0.87$)
in 7 of 9 cases, with the two lowest fits occurring for DeepONet
(Darcy: 0.79, diffusion: 0.84). The lower $R^2$ for DeepONet
may reflect interaction effects between axes---e.g., the branch
network's resolution sensitivity may modulate the effective capacity.

\begin{figure}[t]
\centering
\includegraphics[width=\textwidth]{../figures/fig_3d_exponents.png}
\caption{Summary of 3D scaling exponents $(\alpha, \beta, \gamma)$
across tasks and models. Error bars are 95\% bootstrap CIs.
The data exponent $\alpha$ varies primarily across tasks; the
capacity exponent $\beta$ varies primarily across models; the
resolution exponent $\gamma$ varies across both.}
\label{fig:3d-exponents}
\end{figure}

% ── 5.5: Cross-resolution transfer ──────────────────────────────────
\subsection{Cross-Resolution Transfer}
\label{sec:cross-res}

We evaluate each trained model at $\Reval \in \{32, 64, 128, 256, 512\}$
to construct transfer matrices (\Cref{fig:cross-res}).

\paragraph{Burgers.}
FNO exhibits strong asymmetric transfer: training at $\Rtrain = 256$
yields 12.0\% error at $\Reval = 512$, versus 43.2\% when trained at
$\Rtrain = 32$---a 3.6$\times$ gap. This confirms that FNO retains
high-frequency information from training data and generalizes it to
finer grids. DeepONet transfers more uniformly (error range: 0.37--0.54
across the grid), as its trunk network provides a smooth output basis
that does not depend on training resolution.

\paragraph{Darcy.}
Both models show negligible cross-resolution variation. For FNO, all
entries are within 0.675--0.696 ($<$3\% relative range), indicating
that the Darcy operator is well-represented by low-frequency modes
regardless of grid density. DeepONet exhibits consistently high error
($>1.3$), suggesting it fundamentally struggles with this task.

\paragraph{Diffusion.}
FNO achieves near-perfect resolution invariance: all cross-resolution
entries are $\approx 0.009$ ($<$15\% relative variation). DeepONet
shows mild sensitivity (0.23--0.27), with slightly better performance
when $\Rtrain$ matches $\Reval$ (diagonal entries).

\begin{figure}[t]
\centering
\includegraphics[width=\textwidth]{../figures/fig_cross_res_per_task.png}
\caption{Cross-resolution transfer matrices for DeepONet and FNO
across the three tasks. Each cell shows the mean $\relLtwo$ when
training at $\Rtrain$ (row) and evaluating at $\Reval$ (column),
averaged over all $N$, $D$, and seeds.
Brighter = lower error.}
\label{fig:cross-res}
\end{figure}

% ── 5.6: Best absolute performance ──────────────────────────────────
\subsection{Best Absolute Performance}
\label{sec:best}

\Cref{tab:best} reports the best relative $L^2$ error achieved by each
model on each task, across all hyperparameter configurations.

\begin{table}[h]
\centering
\caption{Best relative $L^2$ error per model per task, and the
configuration that achieves it.}
\label{tab:best}
\begin{tabular}{@{}llrrl@{}}
\toprule
\textbf{Task} & \textbf{Model} & \textbf{$\relLtwo$} & $D$ &
\textbf{Config} \\
\midrule
\multirow{3}{*}{Burgers}
  & MLP      & 0.199 & 263K & $N\!=\!5000, R\!=\!256$, xlarge \\
  & DeepONet & 0.267 & 461K & $N\!=\!5000, R\!=\!256$, xlarge \\
  & FNO      & \textbf{0.048} & 287K & $N\!=\!5000, R\!=\!256$, medium \\
\midrule
\multirow{3}{*}{Darcy}
  & MLP      & \textbf{0.668} & 82K & $N\!=\!5000, R\!=\!32$, large \\
  & DeepONet & 0.734 & 428K & $N\!=\!5000, R\!=\!128$, xlarge \\
  & FNO      & 0.525 & 287K & $N\!=\!1000, R\!=\!64$, medium \\
\midrule
\multirow{3}{*}{Diffusion}
  & MLP      & 0.018 & 82K & $N\!=\!5000, R\!=\!32$, large \\
  & DeepONet & 0.099 & 404K & $N\!=\!500, R\!=\!32$, xlarge \\
  & FNO      & \textbf{0.002} & 124K & $N\!=\!5000, R\!=\!256$, small-med \\
\bottomrule
\end{tabular}
\end{table}

FNO dominates on Burgers (4.8\% vs.\ MLP's 19.9\%) and diffusion
(0.2\% vs.\ MLP's 1.8\%), both temporal evolution tasks for which the
spectral parameterization is well-suited. Darcy is the hardest benchmark
for all models; the best FNO error of 52.5\% suggests that our
1D Darcy formulation poses a challenging approximation problem,
potentially requiring more data or architectural innovations. Notably,
the best MLP configurations for Darcy and diffusion prefer low
resolution ($R = 32$), reflecting that the MLP's parameter count scales
with $R$, making it prone to overfitting at high resolution---a
practical consideration absent from resolution-invariant architectures.


% ══════════════════════════════════════════════════════════════════════
\section{Discussion}
\label{sec:discussion}

\subsection{When Does Resolution Matter?}

Our results identify a clear dichotomy: resolution is a valuable scaling
axis for PDEs with fine-scale spatial structure (advection, shocks) but
provides diminishing returns for smooth operators (diffusion).
Practitioners should:
\begin{itemize}
  \item For advection-dominated problems: invest in higher-resolution
    training data; the resolution exponent $\gamma$ is comparable
    to the data exponent $\alpha$.
  \item For diffusion-dominated problems: allocate budget toward more
    training samples rather than finer grids; even $R = 32$ captures
    the essential dynamics.
  \item For elliptic problems (Darcy): resolution benefits are
    moderate and uncertain (wide CIs); pilot experiments at the target
    resolution should guide allocation.
\end{itemize}

\subsection{Architecture-Specific Insights}

\paragraph{FNO.}
The Fourier Neural Operator consistently achieves the lowest absolute
errors and shows the strongest cross-resolution transfer. Its spectral
inductive bias aligns well with the spectral PDE solvers used for data
generation, potentially overstating its advantage in settings where the
reference solution is computed differently. FNO's low capacity
sensitivity on Burgers ($\beta = 0.57$) but high sensitivity on Darcy
and diffusion ($\beta > 3.7$) suggests the number of retained Fourier
modes is a critical capacity bottleneck.

\paragraph{DeepONet.}
DeepONet shows uniformly low capacity scaling ($\beta < 0.12$) across
all tasks, indicating that its branch--trunk factorization provides a
strong structural prior. However, its absolute errors lag behind FNO on
every task. The high resolution exponent on Burgers ($\gamma = 1.28$)
reveals that the branch network is sensitive to input quality, which may
be ameliorated by improved input encodings (e.g., learned interpolation
rather than our fixed linear interpolation).

\paragraph{MLP Baseline.}
The MLP serves as a useful reference point. Its resolution-tied nature
means $D \propto R$, which creates an implicit coupling between the
resolution and capacity axes. Despite this limitation, it outperforms
DeepONet in absolute error on both Burgers and diffusion, suggesting
that for low-resolution settings, simple architectures remain
competitive.

\subsection{Validity of the Additive Decomposition}

The high $R^2$ values (0.79--0.95) support the additive power-law
assumption across our experimental range. However, several caveats
apply:
\begin{enumerate}
  \item \emph{Bounded $\gamma$ estimates.} Several resolution
    exponents hit the fitting bound ($\gamma = 5$), indicating that
    the resolution axis has saturated and the additive term
    $c\,R^{-\gamma}$ is negligibly small. In such regimes, the
    effective model reduces to $E(N,D) = \Einf + a\,N^{-\alpha} +
    b\,D^{-\beta}$.
  \item \emph{Wide CIs on $\beta$.} Capacity exponents for FNO
    often have wide intervals (e.g., $[0.35, 3.91]$ on Burgers),
    reflecting that 7 capacity levels with 4 resolution slices provide
    limited statistical power for disentangling the capacity signal.
  \item \emph{Interaction effects.} The additive form ignores
    potential $N$--$R$ or $D$--$R$ interactions. Future work may
    explore multiplicative or tensor-product decompositions.
\end{enumerate}

\subsection{Relationship to Paper~1}

This work extends the scaling-law diagnostic framework of
\citet{perrier2025scaling} in three directions: (i)~from
finite-dimensional dynamics to infinite-dimensional operators,
(ii)~from physics-informed loss functions to data-driven operator
learning, and (iii)~from two axes ($N$, $D$) to three axes ($N$, $D$,
$R$). The central finding of Paper~1---that physics priors shift the
error floor rather than the scaling exponent---has an analog here:
resolution-invariant architectures (DeepONet, FNO) do not necessarily
achieve better \emph{scaling exponents} than the MLP baseline, but they
do achieve lower \emph{error floors} on tasks where spectral or mesh-free
representation is advantageous.


% ══════════════════════════════════════════════════════════════════════
\section{Conclusion}
\label{sec:conclusion}

We have presented the first comprehensive scaling study of operator
learning that treats discretization resolution as a first-class variable.
Our three-axis law
$E(N,D,R) = \Einf + a\,N^{-\alpha} + b\,D^{-\beta} + c\,R^{-\gamma}$
fits 18{,}480 experimental data points with $R^2 > 0.79$ across all
model--task combinations and reveals:

\begin{enumerate}
  \item \textbf{Resolution is task-dependent, not
    architecture-dependent.} The resolution exponent $\gamma$ is
    primarily governed by the spatial complexity of the PDE, ranging
    from $\gamma \approx 0.5$ (Burgers, MLP) to $\gamma \to 5$
    (diffusion, all models).
  \item \textbf{Data scaling is the most consistent axis.}
    The data exponent $\alpha$ is well-determined across all
    configurations and varies $2$--$3\times$ across tasks but
    $<2\times$ across models within each task.
  \item \textbf{FNO's spectral bias enables strong cross-resolution
    transfer} on tasks with rich spectral content (Burgers), but
    provides no advantage when the operator is inherently smooth
    (diffusion).
  \item \textbf{DeepONet's branch--trunk factorization is
    capacity-efficient} ($\beta < 0.12$) but resolution-sensitive on
    advection problems ($\gamma = 1.28$ on Burgers).
\end{enumerate}

These findings provide actionable guidance for practitioners: before
scaling resolution, check whether the PDE's spatial complexity warrants
it. For smooth operators, invest in data; for complex operators, invest
in both data and resolution. Cross-resolution transfer matrices offer
a practical diagnostic for assessing whether a model has learned
resolution-transferable representations.

\paragraph{Limitations and future work.}
Our study is limited to 1D PDEs (Burgers, diffusion) and a 1D Darcy
formulation. Extending to 2D and 3D domains---where resolution scaling
is more consequential ($R^d$ cost growth)---is a natural next step. We
also focus on three architectures; incorporating recent advances
(Geo-FNO, GNOT, operator transformers) would broaden the conclusions.
Finally, the additive scaling law may benefit from interaction terms or
multiplicative formulations for settings with strong axis coupling.


% ══════════════════════════════════════════════════════════════════════
\section*{Acknowledgments}
Experiments were conducted on three RunPod GPU instances (NVIDIA A4000,
L4, and RTX 3090). Total experiment wall time was approximately 8 hours
across all pods.


% ══════════════════════════════════════════════════════════════════════
\begin{thebibliography}{20}

\bibitem[Chen and Chen(1995)]{chen1995universal}
T.~Chen and H.~Chen.
\newblock Universal approximation to nonlinear operators by neural networks
  with arbitrary activation functions and its application to dynamical systems.
\newblock \emph{IEEE Trans. on Neural Networks}, 6(4):911--917, 1995.

\bibitem[Hao et~al.(2023)]{hao2023gnot}
Z.~Hao, Z.~Wang, H.~Su, C.~Ying, Y.~Dong, S.~Liu, Z.~Cheng, J.~Song, and
  J.~Zhu.
\newblock {GNOT}: A general neural operator transformer for operator learning.
\newblock In \emph{ICML}, 2023.

\bibitem[Hestness et~al.(2017)]{hestness2017deep}
J.~Hestness, S.~Narang, N.~Ardalani, G.~Diamos, H.~Jun, H.~Kianinejad,
  M.~M.~A. Patwary, Y.~Yang, and Y.~Zhou.
\newblock Deep learning scaling is predictable, empirically.
\newblock \emph{arXiv:1712.00409}, 2017.

\bibitem[Hoffmann et~al.(2022)]{hoffmann2022training}
J.~Hoffmann, S.~Borgeaud, A.~Mensch, E.~Buchatskaya, T.~Cai, E.~Rutherford,
  D.~de~Las~Casas, L.~A. Hendricks, J.~Welbl, A.~Clark, T.~Hennigan,
  E.~Noland, K.~Millican, G.~van~den~Driessche, B.~Damoc, A.~Guy,
  S.~Osindero, K.~Simonyan, E.~Elsen, J.~W. Rae, O.~Vinyals, and
  L.~Sifre.
\newblock Training compute-optimal large language models.
\newblock In \emph{NeurIPS}, 2022.

\bibitem[Kaplan et~al.(2020)]{kaplan2020scaling}
J.~Kaplan, S.~McCandlish, T.~Henighan, T.~P. Brown, B.~Chess, R.~Child,
  S.~Gray, A.~Radford, J.~Wu, and D.~Amodei.
\newblock Scaling laws for neural language models.
\newblock \emph{arXiv:2001.08361}, 2020.

\bibitem[Li et~al.(2021)]{li2021fourier}
Z.~Li, N.~Kovachki, K.~Azizzadenesheli, B.~Liu, K.~Bhatt, A.~Stuart, and
  A.~Anandkumar.
\newblock Fourier neural operator for parametric partial differential equations.
\newblock In \emph{ICLR}, 2021.

\bibitem[Lu et~al.(2021)]{lu2021learning}
L.~Lu, P.~Jin, G.~Pang, Z.~Zhang, and G.~E. Karniadakis.
\newblock Learning nonlinear operators via {DeepONet} based on the universal
  approximation theorem of operators.
\newblock \emph{Nature Machine Intelligence}, 3(3):218--229, 2021.

\bibitem[Perrier(2025)]{perrier2025scaling}
L.~Perrier.
\newblock Scaling-law diagnostics for physics-informed machine learning.
\newblock Working paper, 2025.

\bibitem[Wang et~al.(2021)]{wang2021learning}
S.~Wang, H.~Wang, and P.~Perdikaris.
\newblock Learning the solution operator of parametric partial differential
  equations with physics-informed {DeepONets}.
\newblock \emph{Science Advances}, 7(40), 2021.

\bibitem[Wen et~al.(2022)]{wen2022ufno}
G.~Wen, Z.~Li, K.~Azizzadenesheli, A.~Anandkumar, and S.~M. Benson.
\newblock {U-FNO}---an enhanced {Fourier} neural operator-based deep-learning
  model for multiphase flow.
\newblock \emph{Advances in Water Resources}, 163:104180, 2022.

\end{thebibliography}

% ══════════════════════════════════════════════════════════════════════
\appendix
\section{Experimental Details}
\label{app:details}

\subsection{Capacity Grid}

\begin{table}[h]
\centering
\caption{MLP and DeepONet capacity grid (hidden layers $\times$ hidden
width). FNO uses a separate grid over modes, width, and number of
layers.}
\label{tab:capacity-grid}
\begin{tabular}{@{}lrrr@{}}
\toprule
\textbf{Level} & \textbf{MLP layers} & \textbf{Hidden width} &
\textbf{$\sim D$ (MLP, $R\!=\!128$)} \\
\midrule
tiny      & 2 & 32  & $\sim$5K \\
small     & 2 & 64  & $\sim$13K \\
small-med & 2 & 96  & $\sim$22K \\
medium    & 2 & 128 & $\sim$34K \\
med-large & 2 & 192 & $\sim$62K \\
large     & 2 & 256 & $\sim$82K \\
xlarge    & 3 & 256 & $\sim$148K \\
\bottomrule
\end{tabular}
\end{table}

\subsection{Training Configuration}
All models are trained with Adam (learning rate $10^{-3}$, weight decay
$10^{-5}$), batch size 64, and maximum 5{,}000 epochs with early
stopping (patience 200, monitoring validation relative $L^2$). Data is
split 80/10/10 into train/validation/test. Relative $L^2$ loss
$\|\hat{v} - v\|_2 / \|v\|_2$ is used for training.

\subsection{PDE Solver Details}

\paragraph{Burgers equation.}
$u_t + u u_x = \nu u_{xx}$ on $[0, 2\pi)$ with periodic boundary
conditions, $\nu = 0.01$, $T = 1.0$. Solved with integrating-factor
RK4 in Fourier space with $\frac{2}{3}$-rule dealiasing and
$\Delta t = 10^{-3}$. Random initial conditions are superpositions of
Fourier modes with random phases and amplitudes drawn from a $k^{-2}$
energy spectrum.

\paragraph{Darcy flow.}
$-\nabla \cdot (a(x) \nabla u(x)) = f(x)$ on $[0, 1]$ with Dirichlet
boundary conditions, $f(x) = 1$. The coefficient field $a(x)$ is drawn
from a log-normal random field. Solved via spectral Galerkin method.

\paragraph{Diffusion equation.}
$u_t = \nu u_{xx}$ on $[0, 2\pi)$ with periodic boundary conditions,
$\nu = 0.01$, $T = 1.0$. Exact spectral solution: each Fourier mode
$\hat{u}_k(t) = \hat{u}_k(0) e^{-\nu k^2 t}$. Same initial condition
distribution as Burgers.


\end{document}
